\documentclass[11pt]{article}

% ============================================================
% EXAMEN PARCIAL — UNIDADES 1 Y 2 (evaluación integral)
% Fundamentos y Proceso de Requisitos (U1, 30%) ·
% Elicitación y Análisis de Requisitos (U2, 70%)
% Ingeniería de Requisitos (ISR-401) · UTEQ · 2026-2027 PPA
% ============================================================

\usepackage[utf8]{inputenc}
\usepackage[T1]{fontenc}
% babel no disponible en este entorno
\usepackage[a4paper,top=2cm,bottom=2.2cm,left=1.6cm,right=1.6cm]{geometry}
\usepackage{array}
\usepackage{tabularx}
\usepackage{longtable}
\usepackage{pdflscape}
\usepackage{colortbl}
\usepackage{xcolor}
\usepackage{enumitem}
\usepackage{fancyhdr}
\usepackage{lastpage}
\usepackage[most]{tcolorbox}
\usepackage{booktabs}
\usepackage{amsmath}
\usepackage[hidelinks]{hyperref}

% ---------- Paleta institucional ----------
\definecolor{utegreen}{HTML}{0B7A3B}
\definecolor{utegreendk}{HTML}{075C2C}
\definecolor{midblue}{HTML}{1F5C8B}
\definecolor{gold}{HTML}{D4A017}
\definecolor{rowgray}{HTML}{EFEFEF}
\definecolor{boxgray}{HTML}{F4F6F4}
\definecolor{lightred}{HTML}{FBEAEA}
\definecolor{lightgold}{HTML}{FBF3DC}
\definecolor{lightblue}{HTML}{EAF1F7}
\definecolor{lvlexc}{HTML}{DCEFE3}
\definecolor{lvlsat}{HTML}{EAF1F7}
\definecolor{lvldev}{HTML}{FBF3E2}
\definecolor{lvlins}{HTML}{F7E6E6}

% ---------- Columnas personalizadas ----------
\newcolumntype{L}[1]{>{\raggedright\arraybackslash}p{#1}}
\newcolumntype{C}[1]{>{\centering\arraybackslash}p{#1}}

% ---------- Encabezado / pie ----------
\setlength{\headheight}{30pt}
\pagestyle{fancy}
\fancyhf{}
\lhead{\footnotesize\textbf{Ingeniería de Requisitos (ISR-401)}}
\rhead{\footnotesize UTEQ \,\textcolor{gold}{$\bullet$}\, 2026--2027 PPA}
\lfoot{\footnotesize Examen Parcial --- Unidades 1 y 2 \,$\cdot$\, Uso exclusivo académico}
\rfoot{\footnotesize Página \thepage\ de \pageref{LastPage}}
\renewcommand{\headrulewidth}{0.6pt}
\renewcommand{\footrulewidth}{0.4pt}

\setlength{\parindent}{0pt}
\renewcommand{\arraystretch}{1.25}

% ---------- Atajos ----------
\newcommand{\thdr}[1]{\cellcolor{utegreen}\textcolor{white}{\textbf{#1}}}
\newcommand{\bhdr}[1]{\cellcolor{midblue}\textcolor{white}{\textbf{#1}}}
\newcommand{\pts}[1]{\hfill{\normalsize\textbf{\textcolor{utegreendk}{#1 pts}}}}


\begin{document}

% ============================================================
% PORTADA
% ============================================================
\begin{center}
{\large\textbf{\textcolor{utegreendk}{UNIVERSIDAD TÉCNICA ESTATAL DE QUEVEDO}}}\\[2pt]
{\small Facultad de Ciencias de la Computación \,\textcolor{gold}{$\bullet$}\, Carrera de Software}\\[10pt]
{\LARGE\textbf{\textcolor{utegreendk}{EXAMEN PARCIAL}}}\\[4pt]
{\large\textbf{UNIDADES 1 Y 2 — Evaluación integral}}\\[2pt]
{\normalsize Fundamentos y Proceso de Requisitos \,$\cdot$\, Elicitación y Análisis de Requisitos}
\end{center}

\vspace{4pt}
\noindent\textcolor{gold}{\rule{\linewidth}{1.5pt}}
\vspace{6pt}

\noindent\begin{tabularx}{\linewidth}{|L{0.24\linewidth}|X|}
\hline
\thdr{Asignatura} & Ingeniería de Requisitos (ISR-401) --- 4.º nivel, Carrera de Software\\
\hline
\rowcolor{rowgray}\thdr{Docente} & Guerrero Ulloa Gleiston Cicerón\\
\hline
\thdr{Resultados de aprendizaje evaluados} &
\textbf{RA1 (Unidad 1):} Aplica los fundamentos teóricos de la IR identificando tipos, propiedades y el proceso de requisitos para analizar su aplicabilidad en proyectos de software reales.\newline
\textbf{RA2 (Unidad 2):} Aplica técnicas de elicitación y análisis de requisitos para identificar, clasificar y modelar los requisitos funcionales y no funcionales, asegurando completitud y coherencia.\\
\hline
\rowcolor{rowgray}\thdr{Ponderación por unidad} & Unidad 1: \textbf{3{,}0 puntos} (30\,\%) \quad$\cdot$\quad Unidad 2: \textbf{7{,}0 puntos} (70\,\%)\\
\hline
\thdr{Modalidad / tiempo} & Individual \,$\cdot$\, Presencial \,$\cdot$\, 120 minutos\\
\hline
\rowcolor{rowgray}\thdr{Puntaje total} & 10{,}0 puntos\\
\hline
\thdr{Período} & 2026--2027 PPA\\
\hline
\rowcolor{rowgray}\thdr{Referentes} & ISO/IEC/IEEE 29148:2018; Pohl (2010); ISO/IEC 25010; SWEBOK v4.0 (KA Software Requirements); IREB CPRE FL v3.1.\\
\hline
\end{tabularx}

\vspace{8pt}
\noindent\begin{tabularx}{\linewidth}{|L{0.30\linewidth}|X|}
\hline
\bhdr{Apellidos y Nombres} & \rule{0pt}{14pt}\\
\hline
\bhdr{Cédula / ID} & \rule{0pt}{14pt}\\
\hline
\bhdr{Paralelo \quad / \quad Fecha} & \rule{0pt}{14pt}\\
\hline
\end{tabularx}

% ============================================================
% ALINEACIÓN CURRICULAR
% ============================================================
\section*{\textcolor{utegreendk}{1. ALINEACIÓN CURRICULAR DE LA EVALUACIÓN}}

Esta evaluación es \textbf{integral}: cada tarea se deriva de un único caso de estudio y se alinea con los objetivos específicos de la asignatura, los resultados de aprendizaje de las unidades 1 y 2 y los niveles de la taxonomía de Bloom, conforme al plan analítico vigente.

\vspace{6pt}
\noindent\begin{tabularx}{\linewidth}{|C{0.08\linewidth}|X|C{0.10\linewidth}|C{0.13\linewidth}|L{0.22\linewidth}|C{0.10\linewidth}|}
\hline
\thdr{Tarea} & \thdr{Objetivo específico de la asignatura al que tributa} & \thdr{RA} & \thdr{Unidad} & \thdr{Nivel de Bloom} & \thdr{Puntos}\\
\hline
T1 & OE1: Identificar fundamentos, tipos, clasificaciones y propiedades de los requisitos y su rol en el ciclo de vida. & RA1 & Unidad 1 & Comprender $\rightarrow$ Analizar & 1{,}5\\
\hline
\rowcolor{rowgray}T2 & OE1: Comprender el proceso de requisitos, sus actores y el framework de IR para seleccionar el modelo adecuado. & RA1 & Unidad 1 & Analizar $\rightarrow$ Evaluar & 1{,}5\\
\hline
T3 & OE2: Aplicar técnicas de elicitación para descubrir requisitos funcionales y no funcionales desde diversas fuentes. & RA2 & Unidad 2 & Aplicar $\rightarrow$ Crear & 2{,}0\\
\hline
\rowcolor{rowgray}T4 & OE2 y OE3: Analizar, clasificar, priorizar y negociar requisitos asegurando completitud y coherencia. & RA2 & Unidad 2 & Analizar $\rightarrow$ Evaluar & 2{,}0\\
\hline
T5 & OE3: Modelar conceptualmente los requisitos mediante metas, escenarios, casos de uso y user stories. & RA2 & Unidad 2 & Crear & 2{,}0\\
\hline
\rowcolor{rowgray}T6 & OE4: Formular requisitos no funcionales cuantificados y verificables (ISO/IEC 25010) como base de una ERS de calidad. & RA2 & Unidad 2 & Aplicar & 1{,}0\\
\hline
\multicolumn{5}{|r|}{\cellcolor{lightgold}\textbf{Total}} & \cellcolor{lightgold}\textbf{10{,}0}\\
\hline
\end{tabularx}

% ============================================================
% INSTRUCCIONES Y ADMISIÓN
% ============================================================
\section*{\textcolor{utegreendk}{2. INSTRUCCIONES GENERALES}}

\begin{enumerate}[leftmargin=18pt,itemsep=2pt]
  \item Lea cuidadosamente el \textbf{caso de estudio completo} antes de responder cualquier tarea.
  \item Todas las tareas se derivan del mismo contexto; sus respuestas deben ser \textbf{coherentes entre sí} (mismos actores, requisitos e identificadores).
  \item Fundamente cada respuesta con la \textbf{terminología técnica de la IR} de las unidades 1 y 2.
  \item Las respuestas sin justificación o con términos vagos sin cuantificar («rápido», «seguro», «confiable») no obtendrán el puntaje completo.
  \item No se permite consulta de material externo, dispositivos electrónicos ni comunicación con terceros.
  \item Escriba con letra legible. Si la tarea provee tabla, respétela; si redacta párrafos, organícelos con claridad.
\end{enumerate}

\begin{tcolorbox}[colback=lightred,colframe=red!60!black,
  title=\textbf{Condiciones de admisión (gatekeeper) --- verificación previa a la calificación},
  fonttitle=\bfseries,boxrule=0.8pt,arc=2mm]
El incumplimiento de \textbf{cualquiera} de las siguientes condiciones suspende la revisión con la rúbrica y asigna la \textbf{calificación mínima de 1{,}0/10}. La decisión final de aplicación corresponde al docente y, en el caso 3, al procedimiento disciplinario institucional.
\begin{enumerate}[leftmargin=18pt,itemsep=2pt]
  \item \textbf{Identificación completa} del estudiante (nombres, cédula, paralelo y fecha) en la portada.
  \item \textbf{Pertinencia al caso:} las respuestas se construyen sobre el caso SICAA provisto; respuestas exclusivamente teóricas o de un contexto ajeno no son admisibles.
  \item \textbf{Probidad académica:} evidencia de copia, uso de material no autorizado o suplantación invalida el examen.
\end{enumerate}
\end{tcolorbox}

\newpage
% ============================================================
% CASO DE ESTUDIO
% ============================================================
\section*{\textcolor{utegreendk}{3. CASO DE ESTUDIO: SISTEMA DE GESTIÓN DEL CENTRO DE ACOPIO AGRÍCOLA «SICAA»}}

\subsection*{Contexto general}

La Asociación de Productores «La Esperanza del Río» administra un \textbf{centro de acopio y comercialización de cacao y frutas tropicales} que recibe, en temporada alta, hasta 60 entregas diarias de unos 250 productores asociados de la zona rural. Hoy, la recepción de lotes, el pesaje, el control de calidad y la liquidación de pagos se registran en cuadernos y hojas de cálculo aisladas, lo que provoca demoras en la fila de camiones, reclamos por pagos mal calculados y la imposibilidad de demostrar a los clientes mayoristas el \textbf{origen} de cada lote, condición que estos exigen cada vez más para la exportación.

La directiva ha decidido contratar el desarrollo del \textbf{Sistema de Gestión del Centro de Acopio Agrícola (SICAA)}. Usted ha sido designado \textbf{analista de requisitos} del proyecto. En las primeras visitas al centro de acopio recopiló las siguientes declaraciones:

\begin{tcolorbox}[colback=boxgray,colframe=utegreen,boxrule=0.8pt,arc=2mm]
\textbf{D1 --- Sra. Maldonado, gerente del centro de acopio:}\\
«El sistema debe registrar el pesaje de cada lote que llega y calcular automáticamente el pago al productor según la calidad del producto y el precio del día. Además, el equipo de desarrollo debe entregarnos avances funcionales cada dos semanas y documentar el proyecto según un estándar reconocido, porque la cooperación internacional que financia el proyecto lo audita.»

\vspace{4pt}
\textbf{D2 --- Téc. Rivas, responsable de control de calidad:}\\
«Mi registro tiene que ser rápido. En temporada alta llegan camiones en fila y no puedo demorarme llenando pantallas; si el sistema me atrasa, prefiero seguir con mi cuaderno.»

\vspace{4pt}
\textbf{D3 --- Sr. Cedeño, productor asociado:}\\
«Quiero ver desde mi celular cuánto entregué, qué calidad me asignaron y cuánto me van a pagar. Eso sí: muchos compañeros casi no usan tecnología, así que debe ser fácil de usar y funcionar aunque la señal en el campo sea mala.»

\vspace{4pt}
\textbf{D4 --- Lcda. Zambrano, contadora:}\\
«Los reportes de liquidación deben ser confiables y la información de pagos debe estar segura; nadie que no esté autorizado puede ver ni cambiar los valores. El sistema debe exportar la información al programa contable que ya usamos.»

\vspace{4pt}
\textbf{D5 --- Sr. Lindao, representante del cliente mayorista exportador:}\\
«Para seguir comprando necesito trazabilidad completa: saber de qué finca proviene cada lote, quién lo recibió, qué calidad obtuvo y en qué fecha. Sin eso, no hay certificación de origen y no hay negocio.»

\vspace{4pt}
\textbf{D6 --- Sra. Maldonado, gerente (segunda reunión):}\\
«Ah, y el sistema debe funcionar en las computadoras que ya tenemos en el galpón y respetar la normativa local de protección de datos personales de los asociados.»
\end{tcolorbox}

\subsection*{Lista preliminar de requisitos identificados}

A partir de las reuniones, el equipo elaboró esta lista preliminar (sin depurar):

\noindent\begin{tabularx}{\linewidth}{|C{0.10\linewidth}|X|}
\hline
\thdr{ID} & \thdr{Enunciado preliminar}\\
\hline
R-01 & El sistema registrará el ingreso de cada lote con productor, peso bruto, tara, producto, fecha y hora.\\
\hline
\rowcolor{rowgray}R-02 & El sistema calculará automáticamente la liquidación de pago según la calidad asignada y el precio del día.\\
\hline
R-03 & El productor podrá consultar sus entregas, calidades y pagos desde una aplicación móvil.\\
\hline
\rowcolor{rowgray}R-04 & El sistema generará un código de trazabilidad por lote, vinculado a la finca de origen y al registro de calidad.\\
\hline
R-05 & El sistema enviará una notificación por mensajería a los productores cuando se publique el precio del día.\\
\hline
\rowcolor{rowgray}R-06 & El sistema se integrará con la báscula electrónica del galpón para capturar el peso automáticamente.\\
\hline
R-07 & El usuario podrá personalizar los colores y el tema visual de la aplicación móvil.\\
\hline
\rowcolor{rowgray}R-08 & El sistema exportará los reportes de liquidación al programa contable existente.\\
\hline
\end{tabularx}

\newpage
% ============================================================
% TAREAS — UNIDAD 1
% ============================================================
\section*{\textcolor{utegreendk}{4. TAREAS DE LA UNIDAD 1 \,(3{,}0 puntos)}}

\subsection*{T1. Sistema, contexto y visión \pts{1{,}5}}

\begin{enumerate}[label=\textbf{T1.\arabic*},leftmargin=26pt,itemsep=3pt]
  \item \textbf{(0{,}4)} Clasifique el SICAA según los tipos de sistemas estudiados (información, embebido, software-intensivo, adaptativo) y \textbf{justifique} con dos características del caso. Indique una \textbf{consecuencia} de esa clasificación para la IR del proyecto.
  \item \textbf{(0{,}4)} Defina el \textbf{límite del sistema} y el \textbf{contexto}: mencione dos elementos que quedan \emph{dentro} del sistema, dos que pertenecen al \emph{contexto} y uno que esté en la \emph{frontera} (interfaz), justificando este último.
  \item \textbf{(0{,}3)} Formule la \textbf{visión del sistema} en un único enunciado, alineada con el problema central del caso.
  \item \textbf{(0{,}4)} Seleccione \textbf{tres perspectivas de contexto} (de las ocho estudiadas) pertinentes al SICAA y sustente cada una con evidencia textual del caso.
\end{enumerate}

\subsection*{T2. Tipos de requisitos, proceso y framework de IR \pts{1{,}5}}

\begin{enumerate}[label=\textbf{T2.\arabic*},leftmargin=26pt,itemsep=3pt]
  \item \textbf{(0{,}5)} Clasifique las declaraciones \textbf{D1 a D6} distinguiendo: requisito funcional, requisito no funcional, restricción y \textbf{requisito de proceso} (vs.\ requisito de producto). Una misma declaración puede contener más de un tipo; identifíquelos por separado en una tabla.
  \item \textbf{(0{,}5)} Recomiende el \textbf{modelo de proceso de requisitos} (cascada, iterativo o ágil) más adecuado para el SICAA con \textbf{tres razones contextualizadas} al caso (considere la exigencia de entregas quincenales y la baja familiaridad tecnológica de los usuarios).
  \item \textbf{(0{,}5)} Aplique el \textbf{framework de IR}: identifique los \textbf{actores del proceso} de requisitos del proyecto, señale en qué \textbf{actividad core} se originó la declaración D6, a qué \textbf{tipo de artefacto} (contexto, requisitos o proceso) corresponde cada parte de D6 y qué \textbf{actividad transversal} debe activarse al incorporarla.
\end{enumerate}

% ============================================================
% TAREAS — UNIDAD 2
% ============================================================
\section*{\textcolor{utegreendk}{5. TAREAS DE LA UNIDAD 2 \,(7{,}0 puntos)}}

\subsection*{T3. Planificación de la elicitación \pts{2{,}0}}

\begin{enumerate}[label=\textbf{T3.\arabic*},leftmargin=26pt,itemsep=3pt]
  \item \textbf{(0{,}8)} Construya la \textbf{matriz fuente $\rightarrow$ técnica}: para al menos \textbf{cuatro fuentes de requisitos} del caso (stakeholders, documentos, sistema actual en cuadernos/hojas de cálculo, estándares), seleccione la \textbf{técnica de elicitación} principal más adecuada (debe usar al menos \textbf{tres técnicas distintas}) y justifique cada selección considerando el perfil de la fuente.
  \item \textbf{(0{,}8)} Diseñe el \textbf{guion de entrevista semiestructurada} al Téc.\ Rivas (control de calidad) con \textbf{cinco preguntas}: combine correctamente preguntas abiertas y cerradas, ordénelas con lógica de embudo y oriéntelas a descubrir el flujo real de trabajo y sus restricciones de tiempo.
  \item \textbf{(0{,}4)} Proponga \textbf{una técnica de apoyo} (brainstorming, prototipado, KJ, mapas mentales o checklists) para abordar la resistencia del Téc.\ Rivas («prefiero seguir con mi cuaderno») y describa \textbf{cómo} la aplicaría en una sesión concreta.
\end{enumerate}

\subsection*{T4. Análisis, priorización y negociación \pts{2{,}0}}

\begin{enumerate}[label=\textbf{T4.\arabic*},leftmargin=26pt,itemsep=3pt]
  \item \textbf{(0{,}8)} Priorice los requisitos \textbf{R-01 a R-08} con la técnica \textbf{MoSCoW} para una primera versión del SICAA. Justifique expresamente por qué los clasificados como \emph{Must} son imprescindibles y por qué los \emph{Won't} pueden posponerse.
  \item \textbf{(0{,}6)} Clasifique \textbf{R-01, R-05 y R-07} según el \textbf{modelo de Kano} (básico/obligatorio, de desempeño/unidimensional, atractivo/de deleite) y justifique cada clasificación desde la perspectiva del productor asociado.
  \item \textbf{(0{,}6)} Las declaraciones \textbf{D2 y D5} generan un \textbf{conflicto de requisitos} (registro mínimo y veloz vs.\ trazabilidad completa con más datos por lote). Tipifique el conflicto según las causas estudiadas, y proponga una \textbf{resolución win-win} concreta y técnicamente viable, indicando qué cede y qué gana cada parte.
\end{enumerate}

\subsection*{T5. Metas, escenarios y casos de uso \pts{2{,}0}}

\begin{enumerate}[label=\textbf{T5.\arabic*},leftmargin=26pt,itemsep=3pt]
  \item \textbf{(0{,}6)} Formule la \textbf{meta principal} del SICAA y constrúyale una \textbf{descomposición AND/OR} con al menos \textbf{dos niveles} y \textbf{cinco submetas}, indicando explícitamente qué descomposiciones son AND y cuáles OR.
  \item \textbf{(0{,}4)} Documente \textbf{textualmente} la meta principal aplicando las \textbf{reglas de documentación de metas} estudiadas (identificador, enunciado verificable, justificación y stakeholder responsable, entre otras).
  \item \textbf{(0{,}6)} Especifique el \textbf{caso de uso} «Registrar recepción de lote» en formato detallado: actor primario, precondiciones, \textbf{flujo principal de al menos seis pasos}, \textbf{un flujo alternativo} (p.\ ej., báscula desconectada) y poscondiciones, coherente con R-01, R-04 y R-06.
  \item \textbf{(0{,}4)} Redacte \textbf{una user story} para el Sr.\ Cedeño (productor) derivada de R-03, con el formato canónico y \textbf{tres criterios de aceptación} verificables.
\end{enumerate}

\subsection*{T6. Requisitos no funcionales cuantificados \pts{1{,}0}}

\begin{enumerate}[label=\textbf{T6.\arabic*},leftmargin=26pt,itemsep=3pt]
  \item \textbf{(1{,}0)} Reescriba las necesidades vagas de \textbf{D2} («rápido»), \textbf{D3} («fácil de usar» y «señal mala») y \textbf{D4} («confiables» y «segura») como \textbf{requisitos no funcionales cuantificados}: para cada uno indique la \textbf{característica de calidad ISO/IEC 25010}, la \textbf{métrica}, el \textbf{valor objetivo} y el \textbf{criterio de verificación}. (Mínimo cuatro NFR.)
\end{enumerate}

\newpage
% ============================================================
% RÚBRICA — RESUMEN POR TAREA
% ============================================================
\section*{\textcolor{utegreendk}{6. RÚBRICA DE EVALUACIÓN — CAPACIDADES E INDICADORES DE LOGRO}}

\noindent\begin{tabularx}{\linewidth}{|C{0.07\linewidth}|L{0.28\linewidth}|X|C{0.09\linewidth}|}
\hline
\thdr{Tarea} & \thdr{Capacidad evaluada} & \thdr{Indicadores de logro} & \thdr{Peso}\\
\hline
T1 & Clasifica el sistema, delimita el contexto, formula la visión e identifica perspectivas de contexto. & Tipo de sistema correcto con consecuencia para la IR; límite/contexto/frontera diferenciados con precisión; visión alineada al problema central; perspectivas pertinentes con evidencia textual del caso. & 1{,}5\\
\hline
\rowcolor{rowgray}T2 & Distingue tipos de requisitos (producto vs.\ proceso), selecciona el modelo de proceso y aplica el framework de IR. & Clasificación correcta de D1--D6 incluyendo el requisito de proceso de D1; recomendación con tres razones contextualizadas; actores, actividad core, tipos de artefacto y actividad transversal correctamente identificados. & 1{,}5\\
\hline
T3 & Planifica la elicitación seleccionando técnicas según la fuente y diseña instrumentos. & Matriz fuente$\rightarrow$técnica con $\geq$4 fuentes y $\geq$3 técnicas justificadas por perfil; guion de 5 preguntas con combinación abierta/cerrada y lógica de embudo; técnica de apoyo aplicada a la resistencia del stakeholder con procedimiento concreto. & 2{,}0\\
\hline
\rowcolor{rowgray}T4 & Prioriza con MoSCoW y Kano y resuelve conflictos mediante negociación win-win. & MoSCoW completo (R-01--R-08) con justificación de Must y Won't; Kano correcto y argumentado para R-01, R-05 y R-07; conflicto D2/D5 tipificado y resuelto con propuesta win-win viable y balanceada. & 2{,}0\\
\hline
T5 & Modela metas (AND/OR), escenarios, casos de uso y user stories coherentes con los requisitos. & Descomposición AND/OR con $\geq$2 niveles y $\geq$5 submetas correctamente tipificadas; meta documentada según las reglas textuales; caso de uso detallado completo y coherente con R-01/R-04/R-06; user story canónica con 3 criterios de aceptación verificables. & 2{,}0\\
\hline
\rowcolor{rowgray}T6 & Especifica NFR cuantificados según ISO/IEC 25010. & $\geq$4 NFR con característica ISO 25010 correcta, métrica, valor objetivo y criterio de verificación; eliminación efectiva de la ambigüedad de origen. & 1{,}0\\
\hline
\multicolumn{3}{|r|}{\cellcolor{lightgold}\textbf{Total}} & \cellcolor{lightgold}\textbf{10{,}0}\\
\hline
\end{tabularx}

% ============================================================
% ESCALA Y FÓRMULA
% ============================================================
\section*{\textcolor{utegreendk}{7. ESCALA DE NIVELES Y FÓRMULA DE CALIFICACIÓN}}

\noindent\colorbox{boxgray}{\parbox{\dimexpr\linewidth-2\fboxsep}{\small
\textbf{Escala de niveles.} Cada tarea tiene un \emph{peso} en puntos. El nivel alcanzado determina la fracción del peso obtenida:\\[2pt]
\colorbox{lvlexc}{\textbf{Excelente (N4)} = 100\,\%}\quad
\colorbox{lvlsat}{\textbf{Satisfactorio (N3)} = 75\,\%}\quad
\colorbox{lvldev}{\textbf{En desarrollo (N2)} = 50\,\%}\quad
\colorbox{lvlins}{\textbf{Insuficiente (N1)} = 25\,\%}\quad
\textbf{No respondida} = 0\,\%.\\[4pt]
\textbf{Fórmula.}\;
$\text{Nota} = \sum_{i=1}^{6} \dfrac{\text{Nivel}_i}{4} \times \text{Peso}_i$,\;
con $\text{Nivel}_i \in \{0,1,2,3,4\}$ y $\sum \text{Peso}_i = 10{,}0$.\\[2pt]
\textbf{Piso por gatekeeper:} si se incumple una condición de admisión, $\text{Nota} = 1{,}0$ y no se aplica la rúbrica.
}}

\vspace{6pt}
\begin{center}
\textit{\small La tabla analítica de niveles por tarea continúa en la siguiente página (orientación horizontal).}
\end{center}

% ============================================================
% TABLA ANALÍTICA (LANDSCAPE)
% ============================================================
\begin{landscape}
\thispagestyle{fancy}
\begin{center}
{\large\textbf{\textcolor{utegreendk}{Tabla analítica de niveles de desempeño por tarea}}}
\end{center}
\vspace{2pt}
{\footnotesize
\renewcommand{\arraystretch}{1.18}
\setlength{\tabcolsep}{4pt}

\begin{longtable}{|L{3.2cm}|L{5.6cm}|L{5.3cm}|L{4.9cm}|L{4.3cm}|}
\hline
\rowcolor{utegreen}
\textcolor{white}{\textbf{Tarea (peso)}} &
\cellcolor{utegreendk}\textcolor{white}{\textbf{N4 — Excelente (100\%)}} &
\cellcolor{midblue}\textcolor{white}{\textbf{N3 — Satisfactorio (75\%)}} &
\cellcolor{gold}\textcolor{white}{\textbf{N2 — En desarrollo (50\%)}} &
\cellcolor{red!60!black}\textcolor{white}{\textbf{N1 — Insuficiente (25\%)}} \\
\hline
\endfirsthead
\hline
\rowcolor{utegreen}
\textcolor{white}{\textbf{Tarea (peso)}} &
\cellcolor{utegreendk}\textcolor{white}{\textbf{N4 — Excelente}} &
\cellcolor{midblue}\textcolor{white}{\textbf{N3 — Satisfactorio}} &
\cellcolor{gold}\textcolor{white}{\textbf{N2 — En desarrollo}} &
\cellcolor{red!60!black}\textcolor{white}{\textbf{N1 — Insuficiente}} \\
\hline
\endhead

% ---------- T1 ----------
\cellcolor{lightblue}\textbf{T1. Sistema, contexto y visión}\newline\textit{(1{,}5 — U1)} &
Clasificación correcta y justificada con consecuencia explícita para la IR; dentro/contexto/frontera diferenciados sin errores; visión sintética alineada al problema; tres perspectivas pertinentes con evidencia textual. &
Clasificación y límite correctos con justificaciones parciales; visión válida pero genérica; perspectivas pertinentes con evidencia incompleta. &
Confusión parcial entre límite y contexto, o perspectivas enunciadas sin evidencia del caso; visión que describe funciones en lugar de propósito. &
Clasificación errónea, límite/contexto indistintos o respuestas teóricas sin uso del caso. \\
\hline

% ---------- T2 ----------
\cellcolor{lightblue}\textbf{T2. Tipos de requisitos, proceso y framework}\newline\textit{(1{,}5 — U1)} &
Clasificación completa de D1--D6 que distingue el requisito de \textbf{proceso} (D1); modelo recomendado con tres razones del caso; actores, actividad core, tipos de artefacto y actividad transversal correctos para D6. &
Clasificación mayormente correcta con una omisión (p.\ ej., el requisito de proceso); recomendación válida con razones parcialmente genéricas; framework aplicado con un error menor. &
Clasificación con varios errores u omisión del análisis producto vs.\ proceso; recomendación sin contextualizar; framework identificado de forma incompleta. &
Clasificación incorrecta en su mayoría; sin recomendación fundamentada; framework ausente o erróneo. \\
\hline

% ---------- T3 ----------
\cellcolor{lightblue}\textbf{T3. Planificación de la elicitación}\newline\textit{(2{,}0 — U2)} &
Matriz con $\geq$4 fuentes y $\geq$3 técnicas justificadas según el perfil de cada fuente; guion de 5 preguntas con combinación abierta/cerrada y embudo correcto; técnica de apoyo con procedimiento concreto orientado a la resistencia del stakeholder. &
Matriz completa con justificaciones breves o una técnica subóptima; guion correcto con orden mejorable; técnica de apoyo adecuada pero con aplicación descrita superficialmente. &
Matriz incompleta o con técnicas asignadas sin justificar; preguntas predominantemente cerradas o genéricas; técnica de apoyo solo nombrada. &
Sin matriz o con técnicas inadecuadas a las fuentes; guion ausente o impertinente al caso. \\
\hline

% ---------- T4 ----------
\cellcolor{lightblue}\textbf{T4. Análisis, priorización y negociación}\newline\textit{(2{,}0 — U2)} &
MoSCoW completo y defendible con Must y Won't justificados; Kano correcto y argumentado desde el productor; conflicto D2/D5 tipificado con propuesta win-win viable que explicita cesiones y ganancias de ambas partes. &
MoSCoW completo con 1--2 asignaciones discutibles sin justificar; Kano con un error de categoría; resolución win-win válida pero asimétrica o poco concreta. &
MoSCoW parcial o sin justificación de Must/Won't; Kano confundido en dos casos; conflicto identificado pero resuelto por imposición (gana una sola parte). &
Priorización arbitraria o ausente; Kano incorrecto; conflicto no identificado o sin propuesta de resolución. \\
\hline

% ---------- T5 ----------
\cellcolor{lightblue}\textbf{T5. Metas, escenarios y casos de uso}\newline\textit{(2{,}0 — U2)} &
AND/OR con $\geq$2 niveles y $\geq$5 submetas correctamente tipificadas; meta documentada según las reglas textuales; caso de uso completo (actor, pre, $\geq$6 pasos, alternativo, post) coherente con R-01/R-04/R-06; user story canónica con 3 criterios verificables. &
Modelado completo con errores menores (una relación AND/OR confusa, un paso ambiguo del flujo o un criterio de aceptación no verificable). &
Descomposición plana o sin tipificar AND/OR; caso de uso reducido a lista de funciones; user story sin formato canónico o con criterios genéricos. &
Artefactos ausentes, incoherentes entre sí o sin relación con los requisitos del caso. \\
\hline

% ---------- T6 ----------
\cellcolor{lightblue}\textbf{T6. NFR cuantificados (ISO/IEC 25010)}\newline\textit{(1{,}0 — U2)} &
$\geq$4 NFR con característica correcta, métrica, valor objetivo y criterio de verificación; la ambigüedad de origen queda totalmente eliminada. &
NFR completos con una característica mal asignada o un criterio de verificación débil. &
NFR con métricas pero sin valores objetivo, o con características ISO 25010 confundidas en dos o más casos. &
Reescrituras que mantienen los términos vagos o sin estructura de NFR cuantificado. \\
\hline

\end{longtable}
}
\end{landscape}

% ============================================================
% HOJA DE CÁLCULO
% ============================================================
\section*{\textcolor{utegreendk}{8. HOJA DE CÁLCULO DE LA NOTA}}

\noindent\begin{tabularx}{\linewidth}{|C{0.08\linewidth}|X|C{0.10\linewidth}|C{0.10\linewidth}|C{0.14\linewidth}|C{0.20\linewidth}|}
\hline
\bhdr{Tarea} & \bhdr{Descripción} & \bhdr{Unidad} & \bhdr{Peso} & \bhdr{Nivel (0--4)} & \bhdr{Aporte (Nivel/4 $\times$ Peso)}\\
\hline
T1 & Sistema, contexto y visión & U1 & 1{,}5 & & \\
\hline
\rowcolor{rowgray}T2 & Tipos de requisitos, proceso y framework de IR & U1 & 1{,}5 & & \\
\hline
T3 & Planificación de la elicitación & U2 & 2{,}0 & & \\
\hline
\rowcolor{rowgray}T4 & Análisis, priorización y negociación & U2 & 2{,}0 & & \\
\hline
T5 & Metas, escenarios y casos de uso & U2 & 2{,}0 & & \\
\hline
\rowcolor{rowgray}T6 & NFR cuantificados (ISO/IEC 25010) & U2 & 1{,}0 & & \\
\hline
\multicolumn{3}{|r|}{\cellcolor{lightgold}\textbf{Total}} & \cellcolor{lightgold}\textbf{10{,}0} & &
\cellcolor{lightgold}\rule{0pt}{12pt}\\
\hline
\end{tabularx}

\vspace{6pt}
\noindent\textit{\small Subtotales de control: Unidad 1 = T1 + T2 (máx.\ 3{,}0) \quad$\cdot$\quad Unidad 2 = T3 + T4 + T5 + T6 (máx.\ 7{,}0).}

% ============================================================
% CRITERIOS TRANSVERSALES
% ============================================================
\section*{\textcolor{utegreendk}{9. CRITERIOS TRANSVERSALES DE CALIDAD}}

\begin{tcolorbox}[colback=boxgray,colframe=utegreen,boxrule=0.8pt,arc=2mm]
Aplican a todas las tareas y pueden reducir el nivel asignado dentro de cada una:
\begin{itemize}[leftmargin=18pt,itemsep=2pt]
  \item \textbf{Pertinencia:} las respuestas deben usar el contexto del SICAA, no ejemplos genéricos.
  \item \textbf{Fundamentación:} uso correcto y preciso de la terminología técnica de la IR (unidades 1 y 2).
  \item \textbf{Coherencia:} las respuestas de distintas tareas deben ser consistentes entre sí (mismos actores, requisitos, identificadores y decisiones; p.\ ej., el caso de uso de T5 debe respetar la priorización de T4).
  \item \textbf{Cuantificación:} todo requisito formulado o reformulado debe incluir métricas verificables; los términos vagos sin cuantificar reducen el puntaje proporcionalmente.
  \item \textbf{Extensión:} ni tan breve que carezca de sustento, ni tan extensa que exceda el tiempo disponible.
\end{itemize}
\end{tcolorbox}

% ============================================================
% INTERPRETACIÓN
% ============================================================
\section*{\textcolor{utegreendk}{10. INTERPRETACIÓN DE RESULTADOS}}

\noindent\begin{tabularx}{\linewidth}{|C{0.25\linewidth}|X|}
\hline
\thdr{Banda (nota /10)} & \thdr{Interpretación}\\
\hline
\cellcolor{lvlexc}9{,}0 -- 10{,}0 & Dominio sobresaliente: integra los fundamentos (U1) con la elicitación y el análisis (U2); está en condiciones de liderar la construcción de la ERS del PFC.\\
\hline
\cellcolor{lvlsat}7{,}0 -- 8{,}9 & Dominio satisfactorio: aplica correctamente las técnicas con imprecisiones menores de justificación o cuantificación.\\
\hline
\cellcolor{lvldev}5{,}0 -- 6{,}9 & Dominio parcial: persisten errores conceptuales (clasificación de requisitos, semántica AND/OR, NFR sin cuantificar) que deben corregirse con retroalimentación antes de la Entrega 1B del PFC.\\
\hline
\cellcolor{lvlins}$<$ 5{,}0 & Dominio insuficiente: se recomienda tutoría de refuerzo de los temas 2--7 y reelaboración guiada de los artefactos centrales.\\
\hline
\end{tabularx}

\vspace{14pt}
\noindent\begin{tabularx}{\linewidth}{|X|}
\hline
\cellcolor{lightblue}\textbf{Observaciones / retroalimentación} \\
\hline
\rule{0pt}{42pt} \\
\hline
\end{tabularx}

\vspace{18pt}
\noindent\begin{tabularx}{\linewidth}{C{0.5\linewidth}C{0.5\linewidth}}
\rule{6cm}{0.4pt} & \rule{6cm}{0.4pt} \\
Nota final (sobre 10{,}0) & Firma del docente \\
\end{tabularx}

\vspace{10pt}
\noindent\textcolor{gold}{\rule{\linewidth}{1.5pt}}
\begin{center}
\footnotesize Documento de uso académico --- Ingeniería de Requisitos (ISR-401) --- UTEQ, 2026--2027 PPA.
\end{center}

\end{document}
